% LSPV-4 REPORT::STATE_READBACK1::2026-10-12T17:18:45Z
% !TEX program = xelatex
\documentclass[11pt]{article}
\usepackage[margin=0.88in]{geometry}
\usepackage{amsmath,amssymb,amsthm,booktabs,array,hyperref,xurl,microtype}
\setlength{\emergencystretch}{3em}
\newtheorem{theorem}{Theorem}
\newtheorem{proposition}[theorem]{Proposition}
\newtheorem{definition}[theorem]{Definition}
\newcommand{\SAT}{\operatorname{SAT}}
\newcommand{\rank}{\operatorname{rank}}
\title{STATE\_READBACK1: Data-Dependent Checkpoints\\and Charged Source Access}
\author{LSPV-4 / UNIVERSAL\_CORE\_BRIDGE}
\date{12 October 2026, 17:18 UTC}
\begin{document}
\maketitle
\begin{abstract}
We replace trace-field selection by a joint, explicit checkpoint/access
interface. For any deterministic exact Boolean relation service whose
source-dependent state has at most $2^B$ values and whose later query makes
at most $q$ original-source bit probes, its relation matrix has real rank at
most $2^B\sum_{j\le q}\binom Nj$. The proof is elementary decision-tree
polynomial factorization. A Hamming-code construction disproves the
stronger additive guess $B+q\ge\log_2\rank H$: equality on seven source
bits needs only five retained bits and one later source-bit probe.
We instantiate this construction on 128 actual signed Monotone
3-SAT-$(2,2)$ formulas and verify all 16,384 selected compatibility queries.
The statement constrains a specified relation service; it does not force a
one-bit SAT decider to provide that service. P versus NP remains OPEN.
\end{abstract}

\section{What changes, and what is not claimed}
MODELBRIDGE-0028 gives two equal-length, same-YES clean-core formulas with
opposite residual compatibility. A data-free schedule and the decision bit
therefore cannot reconstruct their compatibility. FIELDMASK1, CAUSALCANCEL1,
and TRACECONJ1 separately show why field masks, live causal paths, and
source-injective traces do not establish solver-paid semantic exposure.
The September 26 continuation-vector report already identifies the central
repair: charge the actual continuation service, including restart and reread.
We adopt that interface requirement rather than claim it as new.

The present advance is a precise quantitative bound for a data-bearing
checkpoint plus its allowed later source access, and a constructive
counterexample to naive bit addition. The scoped mathematical implications
below are PROVED/EXACT by the supplied proofs, with independent review
pending. Finite tests are COMPUTATIONAL. No proof assistant was used. No
unrestricted SAT resource pair, all-input compiler, or separation is claimed.

\section{Exact interface and quantifiers}
Fix $N$, a nonempty source domain $X\subseteq\{0,1\}^N$, a finite query set
$\mathcal A$, and a Boolean relation $H:X\times\mathcal A\to\{0,1\}$.
A deterministic preprocessing procedure reads $x$ before receiving $a$ and
leaves state $s(x)\in\mathcal S$. We require $|\operatorname{im}s|\le2^B$.
The later procedure receives $s(x),a$, and read-only bit access to $x$, and
must return $H(x,a)$ exactly using at most $q$ such bit probes on every legal
pair. It may choose each address adaptively.

Here $B$ measures all retained source-dependent information, including any
relevant program counter, head positions, selected cut identifier, and clock
information. Public length and fixed program text are not source-dependent.
A bound on variable-length payload alone is insufficient if its length or
location carries uncounted information. The state may be an arbitrary
nonlinear function of the source; no special trace representation is imposed.

The quantifier order is
\[
 \forall N,X,\mathcal A,H\ \forall s\ \forall Q:
 \bigl[\forall x\in X\ \forall a\in\mathcal A,
 Q^{x}(s(x),a)=H(x,a),\ \#\mathrm{probes}\le q\bigr].
\]
Preprocessing and query computation are not free when applying this model to
an ordinary machine. The lower bound deliberately ignores their cost, which
only strengthens its validity for the specified service. A source-dependent
whole-formula oracle, an externally supplied source pointer with a powerful
decoder, or an uncounted word read is not a single bit probe.

For SAT, the candidate relation is
\[
 H(F,\alpha)=1\iff\exists w\succeq\alpha\quad w\models F.
\]
It is a two-argument service. Its empty-restriction value is the ordinary
SAT decision, but the ordinary decision task does not require all the other
queries. Direct non-$P$ decision on a fixed NP-complete core would suffice
for $P\ne NP$ without an additional transfer to broader SAT syntax.

\section{A checkpoint-conditioned rank/access theorem}
Put
\[
 D(N,q)=\sum_{j=0}^{\min(q,N)}\binom Nj,
 \qquad H_s=H\big|_{s^{-1}(s)\times\mathcal A}.
\]
\begin{theorem}[STATE-CONDITIONED-RANK-ACCESS1]
For every service satisfying the preceding interface,
\[
 \rank_{\mathbb R}H_s\le D(N,q),\qquad
 \boxed{\rank_{\mathbb R}H\le |\operatorname{im}s|D(N,q)
                         \le2^B D(N,q).}
 \tag{1}
\]
Thus, when $H$ has nonzero rank,
\[
 \log_2\rank H\le B+\log_2 D(N,q).
 \tag{2}
\]
\end{theorem}
\begin{proof}
Fix a state value $s$ and query $a$. The future deterministic source access
is a binary decision tree of depth at most $q$ on its legal state fiber.
Unreachable branches can be filled arbitrarily; no algorithm for discovering
these branches is asserted. Each accepting leaf contributes the product of
at most $q$ tested literals $x_i$ or $1-x_i$. Summing these products and
multilinearizing gives a polynomial $p_{s,a}$ of degree at most $q$ agreeing
with the output on the fiber. Repeated reads do not increase this bound.

Let $V_s$ be the evaluation matrix, on sources in the fiber, of all monomials
of degree at most $q$. Each column of $H_s$ lies in its column span. Hence
$H_s=V_s C_s$ for some coefficient matrix and $\rank H_s\le D(N,q)$.
Stacking the state-fiber row blocks yields
$\rank H\le\sum_s\rank H_s$, which proves (1). Taking logarithms proves (2).
The proof is valid for arbitrary promise subsets and arbitrary state maps.
\end{proof}

At $q=0$, rank is bounded by the number of states. With later access,
query-dependent addresses give a much larger family of possible tests than
one fixed $q$-bit message. This is why replacing $\log D(N,q)$ by $q$ is
invalid. A $w$-bit word probe can be simulated by $w$ bit probes, giving the
conservative replacement $q\mapsto wq$, not the original $q$.

For any fixed $q$, (2) also gives
$B\ge\log_2\rank H-O(q\log N)$. It is an interface tradeoff, not an
exponential-state materialization theorem. Even its strongest source-probe
conclusion cannot require more than reading all $N$ bits when local
computation is unrestricted. Additional computation bounds require their
own argument.

\section{Hamming coding defeats additive state-plus-probe counting}
Let $m=2^r-1$, $r\ge2$. Index the coordinates by all nonzero $r$-bit vectors,
equivalently the integers $1,\ldots,m$. Define
\[
 \sigma(x)=\bigoplus_{i:x_i=1}i,
 \qquad \mathcal C=\{c:\sigma(c)=0\}.
\]
The parity-check matrix has rank $r$, so $|\mathcal C|=2^{m-r}$.
For each $x$, let $c(x)=x$ if $\sigma(x)=0$, and otherwise let
$c(x)=x\oplus e_{\sigma(x)}$. Then $c(x)\in\mathcal C$ and is the unique
codeword within distance one. Uniqueness follows because neither one nor
two distinct nonzero columns sum to zero. This is the usual single-error
Hamming-code geometry \cite{hamming1950}; all needed properties are proved
here rather than imported as an unexplained coding oracle.

Store the $m-r$ systematic information bits of $c(x)$ plus one flag recording
whether $x=c(x)$. Thus
\[
 B=m-r+1.
\]
On a later equality query $y$, compute its corresponding summary.
If summaries differ, answer NO. If both summaries describe the same center,
answer YES. Otherwise $x=c\oplus e_i$ and $y=c\oplus e_j$ are in the same
neighbor sphere. Probe $x_j$ and compare it with $y_j$. The values agree
exactly when $i=j$. This gives exact equality using at most one source-bit
probe.

\begin{proposition}[Failure of additive counting]
For equality on $m=2^r-1$ source bits, the above deterministic service has
$\rank H=2^m$, $q=1$, and $B=m-r+1$. For $r\ge3$,
\[
 B+q=m-r+2<m=\log_2\rank H.
\]
Consequently the generic inequality $B+q\ge\log_2\rank H$ is
FALSIFIED/RETIRED. Inequality (1) remains valid.
\end{proposition}

For $m=7$, there are 16 codewords and 32 retained states, with fibers of size
one or seven. Five stored bits and one probe answer all 16,384 equality
pairs. This does not recover the entire source from six fixed bits: the
address changes with the requested comparison. For $m=15$, $B=12$ and the
same one-probe scheme applies.

\section{An actual signed clean-core realization}
Use the two six-variable templates recovered from MODELBRIDGE-0028:
\begin{align*}
 P={}&\{(1,2,3),(1,2,4),(3,5,6),(4,5,6)\},\\
 Q_0={}&\{(-1,-2,-3),(-1,-2,-4),(-3,-5,-6),(-4,-5,-6)\},\\
 Q_1={}&\{(-1,-2,-4),(-1,-2,-5),(-3,-4,-6),(-3,-5,-6)\},\\
 F_0={}&P\land Q_0,\qquad F_1=P\land Q_1.
\end{align*}
Tuple notation denotes disjunctive clauses. Both templates are proper,
duplicate-free, sign-monochromatic, with two occurrences of each sign per
variable. Assignment integers use variable $i$ at bit $i-1$.
The exhaustive six-variable check gives assignments $a_0=19$ satisfying
only $F_0$, $a_1=39$ satisfying only $F_1$, and $a_*=12$ satisfying both.

For $x\in\{0,1\}^m$, form $F_x$ from $m$ variable-disjoint shifted copies
of $F_{x_j}$. Every $F_x$ is still in signed Monotone 3-SAT-$(2,2)$ and is
satisfiable. Let $\alpha_y$ fully assign block $j$ to $a_{y_j}$. Then
\[
 \boxed{H(F_x,\alpha_y)=1\iff x=y.}
 \tag{3}
\]
Thus the selected compatibility matrix is an identity matrix of rank $2^m$.
This is an easy product subfamily inside the syntax, not an assertion that
this subfamily is itself NP-hard or connected. It is a control against
unqualified per-instance resource transfer.

The source encoding is $\gamma(n)\gamma(k+1)$ followed by $k$ clauses of
three sign-plus-identifier records. Here $n=6m$, $k=8m$, and
$w=\operatorname{bitlength}(6m)$, so
\[
 N=\ell_\gamma(6m)+\ell_\gamma(8m+1)+24m(1+w)=\Theta(m\log m).
\]
Positive blocks precede negative blocks. In the first negative clause of
block $j$, the last identifier is $6j+3$ for template zero and $6j+4$ for
template one. Its least significant bit therefore reveals $1-x_j$ in one
actual source-bit probe at a public computable address.

The implementation validates the whole source in one sequential pass before
creating the sketch. It checks all clause records, not merely an external
template promise; malformed inputs are rejected by this recognizer. A
negative four-clause block and the selector vector are the only temporary
source-dependent construction state needed. Query $y$ names the explicitly
defined $\alpha_y$; this does not implement arbitrary partial-assignment
compatibility. Every encoding bit is a constant, $x_j$, or $1-x_j$; hence
for this subfamily the degree argument may use selector dimension $m$
instead of the weaker encoded dimension $N$.

At $m=7$, the test uses 42 variables, 56 clauses, 1,198 encoded bits, and all
128 formulas. It validates all sources and checks all 16,384 actual CNF
assignment evaluations against the one-probe service. The common witness
$a_*$ verifies that ordinary SAT is YES on every member.

\section{Costs and the unresolved decision bridge}
The code records every source read during preprocessing and every later
source-bit query. Preprocessing reads all $N$ encoded bits, so it is not
free. The Python integer implementation conservatively takes polynomial
bit work: $O(N\log N+m^2)$ preprocessing and $O(m^2)$ query-side bit work,
with $O(m+\log N)$ auxiliary bits in a bit-level implementation of the
stated operations. The test harness itself materializes formulas, expected
answers and matrices and has no small-space claim.

Crucially, $B$ is retained state between stages, not total peak workspace.
The query uses temporary state and reads its $m$-bit context. A constant
number of source probes does not establish constant running time. No
near-linear-time/sublinear-space SAT endpoint is claimed for the prototype.

The actual transfer conditions still to prove for a SAT application are:
\begin{description}
\item[] An actual computation must supply the specified contexts or an exact
cost-preserving reduction to them. A one-bit SAT contract does not do so.
\item[] The retained state must be formed without access to the later context;
otherwise it may simply store $H(x,a)$ in one bit.
\item[] All source-dependent side channels, cut identifiers, input addresses,
and later decoded accesses must be counted in the same execution.
\item[] Any lower bound on the resulting interface must exceed what its
compiler and unrestricted local computation already permit. Rank and source
probes alone do not establish a superpolynomial time lower bound.
\end{description}
This improves a necessary interface theorem without declaring its SAT
realization solved. No universal implementation compiler is demanded.
Fixed code or exponents may be chosen inside a contradiction argument, but
no such selection in this report creates the missing service obligation.

\section{Primary literature and transfer status}
The proof of (1) is a classical decision-tree polynomial argument, developed
explicitly here. Beals et al., Sections 3.1--3.2 and Lemmas 4.1--4.2, are
primary background for query/polynomial modeling; their quantum degree
factor is not substituted for the classical degree-$q$ proof
\cite{beals1998}. G\'al--Miltersen define systematic input-plus-index
structures in Section 1 and prove problem-specific tradeoffs, including
Theorem 2 for prefix sum. Their numerical lower bounds are not imported
into SAT \cite{gal2003}.

A closely matching 2026 framework is Garg et al.'s query-with-sketch model.
Its Lemma 3 reduces a systematic data structure to query-with-sketch, and
Lemma 8 gives a conditional min-entropy route to probe lower bounds
\cite{garg2026}. This supports treating source rereads and retained summaries
jointly. Applying its matrix-powering bounds to SAT would need a new exact
reduction and satisfied entropy hypotheses; no such transfer is claimed.
The final primary CCC PDF was retrieved after the arXiv HTML/PDF fetches
failed. Hamming's construction is a genuine coding-to-service construction,
not a cryptographic hardness assumption.

Our rank theorem is stable under relabeling retained states and source bit
permutations/complements. A fixed oracle whose arguments are constructed
only from state, context and counted probe answers does not add source
information; its real running cost must nevertheless be charged. A
whole-source oracle or low-degree-extension oracle is a different access
primitive and invalidates that comparison. The work is not claimed to
cross relativization, algebrization, or Natural Proofs; see
Aaronson--Wigderson for the precise barrier distinction \cite{aw2009}.

\section{Reproducible tests}
All four invocations use a 25-second alarm and a 30-second shell timeout.
All arithmetic ranks are exact integer elimination, not floating point.
The seed is 202610121659. The finite experiment is not a proof of the
universal rank theorem; that theorem has the symbolic proof above.
\begin{center}
\begin{tabular}{p{0.50\linewidth}rr}
\toprule Audit & Assertions & Seconds\\\midrule
Core CNFs and exhaustive Hamming $m=3,7$ & 67,354 & 0.327\\
Adaptive-tree degree and state-fiber rank & 5,573 & 0.119\\
Hamming $m=15$, all fibers and seeded pairs & 1,064,964 & 3.805\\
Eight easy-family and semantic wrapper controls & 1,058 & 0.040\\\midrule
Total & 1,138,949 & 4.291\\\bottomrule
\end{tabular}
\end{center}
No failure occurred. The rank audit covers 189 parameter cells, 4,968 tree
degree checks, 414 state-fiber ranks, and 189 global ranks. The $m=15$ test
checks all 32,768 summaries and 499,712 pairs, not every pair in the full
$2^{15}\times2^{15}$ domain. The easy controls cover Horn, 2-SAT, affine
parity, treewidth-one, independent products, an UNSAT unit pair, redundancy,
and the signed core. Wrapper tests are semantic identity checks, not timed
simulations of every named wrapper. Two hostile controls invalidate the
query-independent checkpoint and original-bit-access hypotheses.

\section{Corpus coverage, reconciliation and handoff}
Actual main4/main5 entry points, the current root synthesis, AGENTS, CODEX,
and the repository research skill were read. Additional substantive reads
include MODELBRIDGE-0028, FIELDMASK1, CAUSALCANCEL1, TRACECONJ1, and the original
Drive continuation-vector report (modified 27 September). The canonical
MODELBRIDGE-0022 report is already present; no recovery body was invented.
GitHub Hamming, readback, and sketch/rank searches did not surface this
specific Hamming equality construction, but no exhaustive novelty or whole
corpus review is claimed. Personal Context search was discovery only.

The old continuation-service requirement is adopted. Bare trace-field,
causal-path, and schedule criteria remain retired in their proved scopes.
The new theorem constrains an actual data-bearing service without choosing
an invariant internal representation. ARC/WRAPPER semantic exactness and
conditional composition remain credited; no all-input coverage status changes.

The local environment lacks a full repository checkout (git hostname
resolution failed), so no repository-wide preflight or full root build is
claimed. Exact source bodies, script, raw outputs and a narrow LSPV-4
reconciliation patch are preserved separately from publication status. Remote
publication and any blocked root update are recorded in the JSON run log.

\paragraph{Next target.}
Identify a decision-required, context-independent checkpoint or a
cost-preserving decision reduction whose continuation must answer a relation
with a nontrivial rank/access or conditional min-entropy lower bound. First
prove that actual interface connection; do not manufacture a future-query
obligation by changing ordinary SAT into a stronger service. A direct
all-algorithm lower bound for the fixed core remains an alternative.

\begin{thebibliography}{9}
\bibitem{beals1998} R. Beals, H. Buhrman, R. Cleve, M. Mosca and R. de Wolf.
Quantum lower bounds by polynomials. arXiv:quant-ph/9802049v3, 1998.
\url{https://arxiv.org/pdf/quant-ph/9802049}.
\bibitem{gal2003} A. G\'al and P. B. Miltersen.
The cell probe complexity of succinct data structures. BRICS RS-03-44, 2003.
\url{https://www.brics.dk/RS/03/44/BRICS-RS-03-44.pdf}.
\bibitem{hamming1950} R. W. Hamming. Error detecting and error correcting codes.
Bell System Technical Journal 29(2):147--160, 1950.
\url{https://doi.org/10.1002/j.1538-7305.1950.tb00463.x}.
\bibitem{garg2026} S. Garg, S. He, Y. Li, P. A. Papakonstantinou and X. Yang.
Systematic data structure lower bounds via the query-with-sketch model.
CCC 2026, LIPIcs 383, 41:1--41:44. Lemmas 3 and 8.
\url{https://doi.org/10.4230/LIPIcs.CCC.2026.41}.
\bibitem{aw2009} S. Aaronson and A. Wigderson.
Algebrization: A new barrier in complexity theory. ACM TOCT 1(1), 2009.
\url{https://www.scottaaronson.com/papers/alg.pdf}.
\end{thebibliography}
\end{document}
% END_OF_VERIFIED_REPORT::LSPV-4::STATE_READBACK1
